\scanpage{230}
\mspage{176}
% archive sheet number: 229
où j'ai \unsure{même|aussi} marqué en vert les cinq ``\uline{foncteurs} d'\uline{\unsure{inversion}}''
\[
\Phi,\Psi,\Phi_0,\Psi_0,C,
\]
qui fournissent des quasi-inverses des foncteurs sur
les \inserted{catégories} localisées \unsure{déduits|induits} par les \uline{foncteurs d'inclusion}.

\textcircled{3} \uline{Les quatre types de catégories} $\mathrm{HOT}_W$

Sur chacune des catégories de l'étage, il y a
\uline{deux} localiseurs à regarder
\diagram{\begin{tikzcd}[ampersand replacement=\&,row sep=small,column sep=0pt]
 W_S^f \arrow[d,Rightarrow] \& {}=W_S^g=W_S^d=W_S^{\mathrm{univ}}\\
 W_S \&
\end{tikzcd}}
donnant lieu à deux localisés distincts, notés
\[
\mathrm{HOT}_W^{\mathrm{lc}}(S)\qquad\text{et}\qquad\mathrm{HOT}_W^{\mathrm{lc}!}(S),
\]
il se trouve que (grâce à $\Phi_0$ et $\Psi_0$, \deleted{et $C$ pour $\underline{\mathrm{Parf}}_W$}) ces localisés
``sont les \deleted{\unsure{m}}\unsure{même|mêmes}'' pour les 8 catégories de l'étage,
\deleted{\unsure{y inclus}}\inserted{à l'exception de} $\underline{\mathrm{Parf}}_W\,S$ ; plus exactement les foncteurs
d'inclusion \inserted{(à l'exception peut-être de ceux ayant pour source $\underline{\mathrm{Parf}}_W\,S$)}
induisent des équivalences sur les
deux types de localisés, les équivalences quasi-inverses étant obtenues par $\Phi_0$ (aile gauche) ou
$\Psi_0$ (aile droite). \inserted{Quand on localise \uline{pour} $W_S$ (à l'exclusion de $W_S^f$) on trouve aussi \deleted{des} quasi-inverses resp.}
\deleted{\unsure{ou}} par $C\mid\underline{\mathrm{Liss}}_{0W}\,S$
$C\mid\underline{\mathrm{Prop}}_{0W}\,S$, pour les deux inclusions
\diagram{\begin{tikzcd}[ampersand replacement=\&]
\underline{\mathrm{Liss}}_{\unsure{0W|W}}\,S \& \underline{\mathrm{Parf}}_W\,S \arrow[l,hook] \arrow[r,hook,no head] \& \underline{\mathrm{Prop}}_{0W}\,S .\quad {*}
\end{tikzcd}}
% slanted
\margin{cf plus bas ce qu'on peut dire pour les localisés pour $W_S^f$}
Par raison de symétrie, le plus joli semble de
définir
\[
\left\{\begin{aligned}
\mathrm{HOT}_W^{\mathrm{lc}}(S)&=(\underline{\mathrm{Fib}}_W\,S)(W_S^f)^{-1}\\
\mathrm{HOT}_W^{\mathrm{lc}!}(S)&=(\underline{\mathrm{Fib}}_W\,S)(W_S)^{-1},
\end{aligned}\right.
\]
et on a un foncteur de localisation
\[
\mathrm{HOT}_W^{\mathrm{lc}}(S)\xrightarrow{\unsure{\lambda|{?}}}\mathrm{HOT}_W^{\mathrm{lc}!}(S)
\]
induit par l'inclusion
\[
W_S^f\subset W_S .
\]
